\section{Collection management and document lifecycle}
\label{sec:organization}

Collection membership defines the scope available to a query. Parts of
a split document retain the collection assignment, and descriptions
help the agent select documents before opening their bodies. The
operations in Table~\ref{tab:levels} maintain these records as the corpus
changes.

\subsection{Classification, metadata, tags and duplicate review}
\label{sec:org-placement}
\label{sec:org-registries}
\label{sec:org-tags}
\label{sec:org-dedup}

The catalogue returns collections and sub-libraries with descriptions.
Each shelf shows document titles and descriptions within its selected scope.
The ingest agent uses that scope and the parsed content to select a
collection,
then writes descriptions and tags grounded in the body. Duplicate review
presents candidate pairs for inspection. Adjacent parts can describe
the same document while containing distinct source sections.
Figure~\ref{fig:workflow} shows how placement and indexing make the
stored material available to a later query, through catalogue navigation
or directly within a specified scope.

HMRS supports library and sub-library creation, metadata edits, moves,
renames, deletion and forced reindexing. A move updates the document path
and registry entries, while reindexing updates the searchable representation.
\looseness=-1 Integrity checks compare document existence and path across the files on
disk, the \texttt{.tree-fs.json} registry and the library's
\texttt{.knowledge-base.yml}, and inspect index metadata.

Sub-libraries keep their descriptions and indices at their own tree nodes.
When the caller supplies a parent scope, the vector service resolves the
parent and its descendants into the collections searched. The lexical stage
of the two-stage service likewise resolves descendant identifiers. The
evaluation's flat arms (Section~\ref{sec:ladder-opt}) call explicit
collection scopes, which fixes the collections being compared.

\subsection{Lifecycle operations, graph and operational knowledge}
\label{sec:org-lifecycle}
\label{sec:org-graph}
\label{sec:org-experience}

An early rename-and-move trial in a disposable collection exposed stale
chunks under the old document path until a forced reindex. The current rename
operation attempts to remove the old index entries and index the renamed
document; indexing failures are recorded separately from the metadata update.
A post-operation probe checks whether queries resolve to the new path.

The document graph links collections, documents and tags for navigation and
candidate expansion. Operational lessons are also stored as retrievable
documents. Experience tools track reviews, use and freshness and use
credibility tiers to order results. Section~\ref{sec:eval} reports
operational-experience cases alongside the management experiments.

\subsection{Organisation contract}
\label{sec:org-integrity}
\label{sec:org-summary}

The organisation contract links each document to its current owning
collection, descriptive metadata and indexed path. For selected lifecycle
operations, Section~\ref{sec:e2} checks these relationships through file
and metadata reconciliation, description review and retrieval probes.
